\section{Lean}\label{sec:lean-lit}

Les cinq arches de \lp{}, dans lean-pool, sont démontrées en Lean~4 avec Mathlib, sans autre axiome que \texttt{propext}, \texttt{Classical.choice} et \texttt{Quot.sound}. Cette affirmation décrit le développement: les fichiers ne contiennent aucune commande \texttt{\#print axioms}, et nous ne consignons aucune vérification séparée à leur sujet. Pour \texttt{lean/bridges}, l'intégration continue du dépôt de cette édition fait un audit des axiomes à chaque modification: à la version décrite ici, il portait sur 437 déclarations, toutes limitées à \texttt{propext}, \texttt{Classical.choice} et \texttt{Quot.sound}. Les énoncés engendrés sont formalisés, comme corollaires des arches, dans le paquet Lake \texttt{lean/bridges} du dépôt de cette édition (\url{https://github.com/Commutator-IO/grothendieck-project}), qui dépend de lean-pool à une révision fixée. Les fichiers utilisés ici sont les suivants; le paquet en contient trois autres, décrits dans la note associée \cite{bridges-notes}:
\begin{itemize}
\item \lean{Monoids.lean} (\lean{isBooleanTopos_singleObj_iff}, \lean{isDeMorganTopos_singleObj_iff}, \lean{isDeMorganTopos_singleObj_of_monoidWithZero}, \lean{isDeMorganTopos_singleObj_of_commMonoid}, \lean{isDeMorganTopos_singleObj_freeMonoid_unit}, \lean{not_isDeMorganTopos_freeMonoid_two}, \lean{isBooleanTopos_singleObj_of_group}); \lean{Preorders.lean} (\lean{isDeMorganTopos_functor_preorder_iff}, \lean{extremallyDisconnected_iff_confluent}, \lean{moritaEquivalentCat_preorder_iff}); \lean{Combinatorics.lean} (\lean{isDeMorganTopos_presheaf_of_hasInitial}, \lean{not_isDeMorganTopos_graphs}, \lean{graphs_pseudocompl_counterexample}, \lean{not_isDeMorganTopos_simplicialSets}, \lean{isDeMorganTopos_augmentedSimplicialSets}, \lean{isDeMorganTopos_presheaf_fintypeCatSkeleton}, \lean{isDeMorganTopos_treeTopos}, \lean{not_isBooleanTopos_treeTopos}); \lean{EssentialPoints.lean} (\lean{nonempty_essentialPoint_iso_iff}, \lean{exists_iso_essentialPoint}, \lean{essentialPointEndMulEquivOfGroup}, \lean{karoubiRel_iff_rank_eq}, \lean{essentialPointClassesEquivFin}).
\item \lean{QuiversMorita.lean} (Morita et points essentiels des représentations de carquois: \lean{moritaEquivalentCat_paths_iff}, \lean{nonempty_representations_equiv_iff}, \lean{essentialPointPathsEquiv}) et \lean{QuiversCyclic.lean} (les $2^n+1$ topologies du carquois cyclique: \lean{card_grothendieckTopology}).
\item \lean{QuiversDeMorgan.lean} (De Morgan et booléen pour les représentations de carquois: \lean{rightOreCondition_paths_iff}, \lean{isDeMorganTopos_paths_iff}, \lean{isBooleanTopos_paths_iff}; il ne contient aucun énoncé de Morita ni de points essentiels), \lean{Corners.lean} (transfert aux coins), \lean{KaroubiMonoid.lean} (les $n$ classes pour $T_n$) et, pour le journal d'étalonnage, \lean{FiniteMonoids.lean} (unique idéal à droite minimal) et \lean{NatTopologies.lean} (les trois topologies sur $\B\mathbb N$, \lean{nat_card_grothendieckTopology_natCat}, avec les faisceaux de celle du milieu, \lean{isSheaf_finalTopology_iff}).
\end{itemize}
Ne sont pas formalisés: le chapitre d'analyse (Mathlib n'a pas le topos des faisceaux sur un espace avec sa négation), les treillis de sous-espaces au-delà des énoncés de De Morgan et de Morita, les dénombrements pour $PT_n$, $I_n$, $O_n$ et $(\mathbb Z/n,\cdot)$, tous les dénombrements de sous-topos de monoïdes, et la description des faisceaux du carquois cyclique. Les énumérations par ordinateur pour $(\mathbb Z/n,\cdot)$ (\S\ref{sec:linear}) ne font pas partie du développement Lean; pour $n>30$, elles passent par la classification des idéaux idempotents et ne vérifient donc pas les topologies indépendamment.
Dans Mathlib, \texttt{SingleObj M} compose selon $f\gg g = g\cdot f$; avec cette convention, la condition d'Ore sur \texttt{SingleObj M} s'écrit $\forall a\,b,\ \exists c\,d,\ a\cdot c = b\cdot d$ (\lean{rightOreCondition_singleObj_iff}), comme dans le texte.

\begin{center}\scriptsize\emergencystretch=3em
\begingroup\def\UrlBreaks{\do\_\do\.\do\/}
\begin{tabularx}{\linewidth}{>{\raggedright\arraybackslash}p{2.6cm}>{\raggedright\arraybackslash}X>{\raggedright\arraybackslash}p{2.9cm}>{\raggedright\arraybackslash}p{2.2cm}}
\toprule
énoncé & déclaration Lean & fichier & statut\\
\midrule
$M$-ensembles booléens si et seulement si groupe; De Morgan si et seulement si Ore & \lean{isBooleanTopos_singleObj_iff}, \lean{isDeMorganTopos_singleObj_iff} & \lean{Monoids} & démontré en Lean\\
carquois: De Morgan, booléen (th.~\ref{th:quivers}) & \lean{isDeMorganTopos_paths_iff}, \lean{isBooleanTopos_paths_iff} & \lean{QuiversDeMorgan} & démontré en Lean (convention de Lean)\\
automates sur $\Sigma$ & \lean{isDeMorganTopos_singleObj_freeMonoid_unit}, \lean{not_isDeMorganTopos_freeMonoid_two} & \lean{Monoids} & démontré en Lean\\
carquois: Morita, points essentiels & \lean{nonempty_representations_equiv_iff}, \lean{essentialPointPathsEquiv} & \lean{QuiversMorita} & démontré en Lean\\
carquois cyclique: $2^n+1$ topologies & \lean{card_grothendieckTopology} & \lean{QuiversCyclic} & démontré en Lean pour le nombre de topologies; les faisceaux de chacune, à la main seulement\\
$T_n$: $n$ points essentiels & \lean{karoubiClassesEquivFin} & \lean{KaroubiMonoid} & démontré en Lean (classes de l'enveloppe de Karoubi)\\
$M_n(\mathbb F_q)$: $n+1$ points essentiels & \lean{essentialPointClassesEquivFin} & \lean{EssentialPoints} & démontré en Lean, sur un corps quelconque\\
points essentiels de $PT_n$, $I_n$, $O_n$; tous les dénombrements de sous-topos du \S\ref{sec:finite} & --- & --- & à la main seulement\\
\bottomrule
\end{tabularx}

\begin{tabularx}{\linewidth}{>{\raggedright\arraybackslash}p{2.6cm}>{\raggedright\arraybackslash}X>{\raggedright\arraybackslash}p{2.9cm}>{\raggedright\arraybackslash}p{2.2cm}}
\toprule
énoncé & déclaration Lean & fichier & statut\\
\midrule
coins: Ore se transfère & \lean{forall_exists_mul_eq_mul_iff_cornerMonoid} & \lean{Corners} & démontré en Lean\\
$M_n(k)$-ensembles, $\mathsf{Vect}_{\le n}$ & --- & --- & à la main seulement (les classes de rang sont en Lean, voir plus haut)\\
treillis de sous-espaces: De Morgan, Morita & \lean{isDeMorganTopos_functor_preorder_iff}, \lean{moritaEquivalentCat_preorder_iff} & \lean{Preorders} & démontré en Lean; le reste à la main seulement\\
dénombrements pour $(\mathbb Z/n,\cdot)$ & --- & --- & énumérés par ordinateur (pas en Lean); preuve à la main seulement\\
$X$ compact séparé & --- & --- & à la main seulement\\
graphes, ensembles simpliciaux, augmentés, $\mathsf{Fin}$, arbres & \lean{not_isDeMorganTopos_graphs}, \lean{not_isDeMorganTopos_simplicialSets}, \lean{isDeMorganTopos_augmentedSimplicialSets}, \lean{isDeMorganTopos_presheaf_fintypeCatSkeleton}, \lean{isDeMorganTopos_treeTopos} & \lean{Combinatorics} & démontré en Lean; les autres lignes combinatoires à la main seulement\\
\bottomrule
\end{tabularx}
\endgroup
\end{center}
